\documentclass[11pt,a4paper]{article}

\usepackage[margin=1in]{geometry}
\usepackage{amsmath,amssymb,amsthm,mathtools}
\usepackage{microtype}
\emergencystretch=3em
\usepackage[T1]{fontenc}
\usepackage[utf8]{inputenc}
\usepackage{lmodern}
\usepackage{booktabs}
\usepackage{enumitem}
\usepackage[dvipsnames]{xcolor}
\usepackage{hyperref}
\usepackage[nameinlink,capitalize]{cleveref}

\hypersetup{
  colorlinks=true,
  linkcolor=blue!50!black,
  citecolor=blue!50!black,
  urlcolor=blue!50!black,
  pdftitle={Inverting a 2D LEEM focal series: the fastest still-correct object reconstruction},
  pdfauthor={Wilson Lau}
}

\newtheorem{theorem}{Theorem}[section]
\newtheorem{lemma}[theorem]{Lemma}
\newtheorem{proposition}[theorem]{Proposition}
\newtheorem{corollary}[theorem]{Corollary}
\theoremstyle{definition}
\newtheorem{definition}[theorem]{Definition}
\theoremstyle{remark}
\newtheorem{remark}[theorem]{Remark}

\newcommand{\R}{\mathbb{R}}
\newcommand{\C}{\mathbb{C}}
\newcommand{\Z}{\mathbb{Z}}
\newcommand{\dd}{\mathrm{d}}
\newcommand{\ee}{\mathrm{e}}
\newcommand{\ii}{\mathrm{i}}
\newcommand{\conj}[1]{\overline{#1}}
\newcommand{\abs}[1]{\lvert#1\rvert}
\newcommand{\norm}[1]{\lVert#1\rVert}
\newcommand{\qap}{q_{\mathrm{ap}}}
\newcommand{\Dz}{\Delta z}
\newcommand{\RFO}{R_{\mathrm{FO}}}
\newcommand{\RCTF}{R_{\mathrm{CTF}}}
\newcommand{\Nap}{N_{\mathrm{ap}}}
\newcommand{\Nf}{N_{\mathrm{f}}}

\title{Inverting a 2D LEEM focal series:\\
the fastest still-correct object reconstruction}
\author{Wilson Lau}
\date{24 August 2026}

\begin{document}
\maketitle

\begin{abstract}
Given an experimental through-focal stack of 2D LEEM intensities
$I(\mathbf{r},\Dz_k)$, recover the object exit wave
$\psi_0(\mathbf{r})=\sigma(\mathbf{r})\,\ee^{\ii\phi(\mathbf{r})}$.
Yu, Lau and Altman~\cite{Yu2019} never invert: they forward-simulate
the bilinear Fourier-optics (FO) kernel $\RFO$ and match contrast.
Five inverses were compared. The unique published 2D LEEM exit-wave
reconstruction is linear Schiske/Wiener restoration~\cite{Duden2014},
which is $O(K N\log N)$ on $N$ pixels and $K$ defoci, fills contrast-transfer
(CTF) zeros by stacking defoci, and is the Fr\'echet derivative of
bilinear FO on the axes $\mathbf{q}'=\mathbf{0}$.
It is biased for the strong-phase objects (graphene ripples, atomic steps)
that make FO necessary in the first place.
The recommended pipeline is therefore two-stage:
\textbf{(1)~Fourier-diagonal multi-defocus Tikhonov} on the CTF slice
$\RFO(\mathbf{q},0,\Dz)$ --- the asymptotically fastest method that still
fills CTF zeros --- then \textbf{(2)~optional Gauss--Newton} on the full
bilinear FO residual, skipped when that residual already sits at the
noise floor ($\max\abs{\phi}\lesssim 0.3$).
For Yu's 1D sinusoid the second stage collapses to Jacobi--Anger
fitting of $M=2\lfloor \qap\Lambda\rfloor+1$ Bessel modes
($O(M^2)$ per residual), which is the fastest still-correct method in
that class.
Lean~4 proves uniqueness, the bias--noise identity, the triangle error
bound, aperture invisibility, gauge invariance, the quadratic remainder,
and an $O(KN\log N)$ cost model, without claiming FFT existence.
\end{abstract}

\tableofcontents
\newpage

\section{The inverse problem}
\label{sec:problem}

A LEEM experiment records a real, nonnegative intensity
$I(\mathbf{r})$ in a $M=1$ image plane, typically as a through-focal
series $I(\mathbf{r},\Dz_k)$, $k=1,\dots,K$.
The object is the complex exit wave
\begin{equation}
\psi_0(\mathbf{r})=\sigma(\mathbf{r})\,\ee^{\ii\phi(\mathbf{r})},
\qquad
\boldsymbol{\Psi}=\mathcal{F}[\psi_0],
\end{equation}
with $\phi$ in radians.
Yu, Lau and Altman~\cite{Yu2019} write the forward map as the bilinear
form
\begin{equation}
\label{eq:I}
I(\mathbf{r},\Dz)
=\iint
\boldsymbol{\Psi}(\mathbf{q})\,\boldsymbol{\Psi}^*(\mathbf{q}')
\,\RFO(\mathbf{q},\mathbf{q}',\Dz)
\,\ee^{2\pi\ii(\mathbf{q}-\mathbf{q}')\cdot\mathbf{r}}
\,\dd^2q\,\dd^2q',
\end{equation}
where $\RFO=R_0 E_S E_{C,\mathrm{tot}}$ is supported on the contrast
aperture $\abs{\mathbf{q}},\abs{\mathbf{q}'}\le\qap$.
The CTF kernel $\RCTF$ agrees with $\RFO$ on the axes
$\mathbf{q}'=\mathbf{0}$ (Lean: \texttt{R\_FO\_eq\_R\_CTF\_axis}).

\paragraph{What is not recoverable.}
\begin{itemize}[leftmargin=*,itemsep=0.2em]
\item Global phase: $I(\ee^{\ii\theta}\psi_0)=I(\psi_0)$
(Lean: \texttt{ihat\_gauge}, \texttt{objectWave\_phase\_shift}).
\item Modes $\abs{\mathbf{q}}>\qap$: $\RFO(\mathbf{q},0,\Dz)=0$ for every
defocus, so the Tikhonov estimator of that bin is identically $0$
(\texttt{R\_FO\_axis\_eq\_zero\_of\_outside},
\texttt{tikhonovXhat\_outside\_aperture}).
\item A single image of a general complex object: one complex linear
measurement of the pair $\bigl(X(\mathbf{q}),X(-\mathbf{q})\bigr)$ always
has a kernel (\texttt{one\_measurement\_not\_injective}).
Need $K\ge 2$ defoci with independent CTF slices.
\end{itemize}

Yu~\emph{et al.}~\cite{Yu2017,Yu2019} \emph{do not invert}.
They choose a trial object (step, sinusoid), run~\eqref{eq:I}, and
overlay contrast versus $\Dz$ on an experimental series.
The only demonstrated 2D experimental LEEM exit-wave reconstruction is
Duden, Thust, Kumpf and Tautz~\cite{Duden2014}: linear Schiske restoration
of a molecular monolayer, assuming a strong specular beam.

\section{Five candidates}
\label{sec:five}

Write $N$ for the number of image pixels, $K$ for the number of
defoci, and $\Nap$ for the number of Fourier samples inside the
aperture ($\Nap\sim \pi(\qap L)^2$ in 2D).
Typical LEEM numbers: $N=512^2$, $K=11$--$21$,
$\Nap\sim 10^3$--$10^5$.

\subsection*{1. Single-image Wiener / CTF deconvolution}
One FFT, divide by $\RCTF(\mathbf{q},\Dz)$, Tikhonov at the zeros.
Cost $O(N\log N)$.
Scientifically inadequate: CTF zeros punch holes in the spectrum, and
a general complex object is not identified by one plane
(\cref{sec:problem}).
Fastest, and the wrong inverse.

\subsection*{2. Transport of intensity (TIE)}
Two or three closely spaced defoci, FFT Poisson solve, cost
$O(N\log N)$~\cite{Teague1983,Paganin}.
Assumes paraxial coherent $\partial_z I=-(\lambda/2\pi)\nabla\cdot(I\nabla\phi)$.
LEEM after $15\,\mathrm{kV}$ acceleration has $\lambda\sim 10\,\mathrm{pm}$,
a hard aperture, $C_3/C_5$, and partial coherence.
TIE ignores all of that; it is a low-pass phase map, not an FO inverse.
Not used on experimental LEEM graphene/steps.

\subsection*{3. Multi-defocus Fourier-diagonal Tikhonov (Schiske)}
\label{sec:schiske}
After $K$ FFTs, each spatial frequency $\mathbf{q}$ decouples.
The linearized observations are
\begin{equation}
\label{eq:linear}
\hat I_k(\mathbf{q})
\approx
h_k(\mathbf{q})\,X(\mathbf{q})
+
g_k(\mathbf{q})\,\overline{X(-\mathbf{q})},
\qquad
h_k=\RFO(\mathbf{q},0,\Dz_k).
\end{equation}
The scalar Tikhonov estimator of a one-branch unknown
(Lean: \texttt{tikhonovXhat}) is
\begin{equation}
\label{eq:xhat}
\hat x
=
\frac{\sum_k \overline{h_k}\,y_k}{\alpha+\sum_k\abs{h_k}^2}.
\end{equation}
Cost: $K$ forward FFTs, $O(KN)$ bin arithmetic, one inverse FFT
--- $O(KN\log N)$.
This is Duden's algorithm with the FO/CTF envelopes of this library
in place of a coherent probe.
It is the Fr\'echet derivative of~\eqref{eq:I} at vacuum
(\texttt{ihatJac\_vacuum}): one Gauss--Newton step from vacuum
\emph{is}~\eqref{eq:xhat}.

\subsection*{4. Gerchberg--Saxton / HIO / IWFR}
Iterate modulus constraints between planes linked by a propagator.
Cost $O(T K N\log N)$ with $T\sim 10$--$10^3$ iterations.
Under partial coherence the intensity constraint is the quadratic set
of~\eqref{eq:I}; the metric projector onto that set \emph{is} a
Gauss--Newton step.
Naive modulus replacement is the coherent CTF model and is biased
when $\Gamma_C\Gamma_S\neq 1$ off axis~\cite{Yu2017}.
Either incorrect, or a slower undamped rewrite of method~5.

\subsection*{5. Nonlinear least squares / Gauss--Newton on bilinear FO}
Minimize $\sum_k\norm{I_{\mathrm{FO}}(\psi;\Dz_k)-I_k^{\mathrm{obs}}}^2$.
Each residual is a difference-frequency accumulation of the TCC,
$O(K\Nap^2+KN\log N)$ per iteration, $T\sim 5$--$20$ if initialized
by method~3.
This is the inverse of the kernel Yu actually uses, and the right
estimator for strong-phase objects.
It is not Fourier-diagonal about a generic background
(\texttt{ihat\_add}), so it is not $O(KN\log N)$.

\section{Verdict}
\label{sec:verdict}

\begin{center}
\small
\begin{tabular}{@{}lccc@{}}
\toprule
Method & Cost & Fills CTF zeros & FO-faithful for $\phi\not\ll 1$ \\
\midrule
1. Single-image Wiener & $O(N\log N)$ & no & no \\
2. TIE & $O(N\log N)$ & n/a (wrong PDE) & no \\
\textbf{3. Multi-defocus Tikhonov} & $\boldsymbol{O(KN\log N)}$ & \textbf{yes} ($K\ge 2$) & linearized only \\
4. GS / HIO & $O(TKN\log N)$ & if $K\ge 2$ & only if the projector is FO \\
5. Gauss--Newton on FO & $O(TK\Nap^2)$ & yes & \textbf{yes} \\
Jacobi--Anger LS (1D sinusoid) & $O(M^2)$ & n/a (parametric) & \textbf{yes} \\
\bottomrule
\end{tabular}
\end{center}

\paragraph{Fastest usable method for a general 2D experimental stack.}
Method~3: multi-defocus Fourier-diagonal Tikhonov on the CTF slice.
It is the cheapest algorithm that (i)~reads all $KN$ Fourier samples,
(ii)~fills interior CTF zeros of $\sin(2\pi\chi_S)$ by stacking defoci,
and (iii)~has a unique minimizer for every $\alpha>0$.
Any method that uses the whole stack is $\Omega(KN)$ after the FFTs;
dense bilinear apply is $\Theta(KN^2)$ and is already strictly slower
at $N=128$ (Lean: \texttt{reconstructCost\_lt\_dense\_128}).

\paragraph{Best FO-faithful method.}
Method~3 as initializer, then a few damped Gauss--Newton steps on
method~5 (method~D / MAL~\cite{Coene1996} if aberrations are unknown).
Skip Gauss--Newton when the FO residual of method~3 is at the noise
floor --- equivalently $\max\abs{\phi}\lesssim 0.3$ at typical LEEM SNR,
i.e.\ topography $A\lesssim 0.01\,\mathrm{nm}$ at $E_0=10\,\mathrm{eV}$.

\paragraph{Fastest still-correct method for Yu's sinusoid.}
Jacobi--Anger fitting of $\varphi$ on $M=2\lfloor\qap\Lambda\rfloor+1$
modes, \cref{sec:bessel}.
This is method~5 restricted to a one-parameter manifold, cost
$O(M^2)$ with $M\ll\Nap$.

\section{Stage 1: the discrete algorithm}
\label{sec:algo}

\begin{enumerate}[leftmargin=*,itemsep=0.35em]
\item Align the stack $I_k(\mathbf{r})$ (rigid shift; not formalized).
\item Discrete Fourier transform each plane,
$\hat I_k=\mathrm{DFT}[I_k-\langle I_k\rangle]$ (mean subtracted so
the DC bin is handled separately: $\RFO(0,0)=1$ recovers total flux).
\item For every frequency $\mathbf{q}$ with $\abs{\mathbf{q}}\le\qap$
and every defocus $k$,
$h_k(\mathbf{q})=\RFO(\mathbf{q},0,\Dz_k)$
from the closed-form envelopes already proved in this library.
\item Scalar (or $2\times 2$ conjugate-pair) Tikhonov~\eqref{eq:xhat}
with $\alpha>0$.
\item Set $\hat X(\mathbf{q})=0$ for $\abs{\mathbf{q}}>\qap$.
\item Inverse DFT $\hat\psi=\mathrm{iDFT}[\hat X]$, then
$\sigma=\abs{\hat\psi}$, $\phi=\arg(\hat\psi)$ in the vacuum gauge
$\hat X(\mathbf{0})$ real and positive.
\end{enumerate}

\paragraph{Units.}
$\RFO$ uses column wavelength $\lambda$ (after $U_a$), not $\lambda_0$.
Object phase uses $\lambda_0$: $\phi=4\pi A/\lambda_0$ for a topographic
height $A$ in reflection.

\paragraph{Defocus sign.}
This library uses $+\tfrac12\Dz\lambda q^2$ in $\chi_S$.
Experimental lens current must be mapped onto that convention by one
fitted affine transform; mixing the Yu GitHub minus-sign convention
reverses overfocus and underfocus.

\section{Error}
\label{sec:error}

All statements below are finite-dimensional, with defocus index in a
finite type $\kappa$ (in Lean: any \texttt{Fintype}).
Noise $n_k$ is a deterministic discrepancy (linearization remainder $+$
instrument noise $+$ discretization), not a random variable.

\begin{theorem}[Unique minimizer]
\label{thm:unique}
For $\alpha>0$ the energy
$J(x)=\sum_k\norm{h_k x-y_k}^2+\alpha\norm{x}^2$
has a unique global minimizer~\eqref{eq:xhat}.
Completing the square:
$J(x)=J(\hat x)+D\norm{x-\hat x}^2$ with $D=\alpha+\sum\abs{h_k}^2>0$.
\end{theorem}
Lean: \texttt{tikhonovJ\_min}, \texttt{tikhonov\_unique},
\texttt{tikhonovJ\_eq\_shift}.
The same identity gives uniqueness at $\alpha=0$ whenever
$\sum\abs{h_k}^2>0$.

\begin{theorem}[Bias--noise identity]
\label{thm:error}
If $y_k=h_k x^\star+n_k$ and $D\neq 0$, then
\begin{equation}
\hat x-x^\star
=
\frac{\sum_k\overline{h_k}\,n_k-\alpha x^\star}{D}.
\end{equation}
\end{theorem}
Lean: \texttt{tikhonov\_error}.
Taking $\mathbb{E}[n]=0$ is informal (no probability space).

\begin{theorem}[Triangle bound, sharp]
\label{thm:bound}
If $\alpha\ge 0$ and $D>0$,
\begin{equation}
\abs{\hat x-x^\star}
\le
\frac{\sum_k\abs{h_k}\,\abs{n_k}+\alpha\,\abs{x^\star}}{D}.
\end{equation}
Equality holds for one measurement, $h=n=1$, $x^\star=-1$, $\alpha=1$.
\end{theorem}
Lean: \texttt{tikhonov\_error\_bound},
\texttt{tikhonov\_error\_bound\_sharp}.

\begin{theorem}[Quadratic remainder of bilinear FO]
\label{thm:quad}
With vacuum spectrum $1_{\mathbf{q}=\mathbf{0}}$ and increment $\delta$,
\begin{equation}
\hat I(\mathrm{vac}+\delta)
=
\hat I(\mathrm{vac})
+DI_{\mathrm{vac}}[\delta]
+\hat I(\delta),
\end{equation}
where $DI_{\mathrm{vac}}[\delta](\xi)=\RFO(\xi,0)\,\delta(\xi)+\RFO(0,-\xi)\,\overline{\delta(-\xi)}$
and $\abs{\hat I(\delta)(\xi)}\le\bigl(\sum\abs{\RFO}\bigr)\bigl(\sum\abs{\delta}\bigr)^2$.
\end{theorem}
Lean: \texttt{ihat\_add}, \texttt{ihatJac\_vacuum},
\texttt{ihat\_quadratic\_remainder}, \texttt{ihat\_bound}.
The dropped term is exactly bilinear: linearized inversion is biased by
$O(\norm{\delta}^2)$, which is the whole image for a strong-phase
sinusoid~\cite{Yu2017}.

\begin{theorem}[Interior CTF zeros of a NAC column]
\label{thm:ctfzero}
$\sin(2\pi\chi_S)=0$ whenever $2\chi_S\in\Z$.
For $C_3=C_5=0$, $\Dz\neq 0$, and a large enough aperture, there exists
nonzero $q$ inside the disk with that vanishing
(\texttt{exists\_interior\_weakPhase\_zero}).
A single defocus is therefore blind at those rings; $K\ge 2$ with
incommensurate $\Dz$ typically fills them.
\end{theorem}

The $2\times 2$ Gram determinant of the vacuum Jacobian is positive for
$\alpha>0$ (\texttt{gramDet\_pos}), so the conjugate-pair normal
equations are invertible even though one measurement is not
(\cref{sec:problem}).

\section{Complexity}
\label{sec:cost}

No FFT existence theorem is claimed.
The cost model (Lean: \texttt{dftCost}, \texttt{reconstructCost}) is
\begin{equation}
\texttt{dftCost}(N):=N\log_2 N,
\qquad
\texttt{reconstructCost}(K,N)
=(K+1)\,N\log_2 N+N(8K+4).
\end{equation}
The $8K+4$ term is a scalar (or $2\times 2$) Gram/right-hand side.

\begin{theorem}[Modelled cost]
\label{thm:cost}
For $K\ge 1$ and $N\ge 2$,
$\texttt{reconstructCost}(K,N)\le 14\,K\,N\log_2 N$.
At $N=128$ this is strictly cheaper than a dense bilinear apply
$KN^2$, for every $K\ge 1$.
\end{theorem}
Lean: \texttt{reconstructCost\_le},
\texttt{reconstructCost\_lt\_dense\_128}.
The constant $14$ is attained in the model at $K=1$, $N=2$.

\paragraph{Wall-clock order of magnitude.}
For $n=512$ ($N=n^2\approx 2.6\times 10^5$), $K=21$:
about $22$ two-dimensional FFTs plus $O(KN)$ arithmetic ---
milliseconds to a few tens of milliseconds on a laptop.
Gauss--Newton at $\Nap\sim 4\times 10^3$ is $\sim 16\times 10^6$ pair
evaluations per residual per defocus: seconds, not milliseconds.
Gerchberg--Saxton at $T=100$ is two orders of magnitude above stage~1
and does not improve the FO model.

\section{Stage 2 and the 1D sinusoid}
\label{sec:bessel}

Kinematic LEEM phase on reflection is $\varphi=4\pi A/\lambda_0$.
At $E_0=10\,\mathrm{eV}$, $\lambda_0\simeq 0.39\,\mathrm{nm}$, so
$A=0.1\,\mathrm{nm}$ already gives $\varphi\sim 3$: not weak.
Jacobi--Anger (Lean: \texttt{jacobi\_anger})
\begin{equation}
\ee^{\ii\varphi\sin(2\pi x/\Lambda)}
=\sum_{n\in\Z}J_n(\varphi)\,\ee^{2\pi\ii n x/\Lambda}
\end{equation}
places intensity on off-axis pairs $(n,m)$ where $\Gamma_C\Gamma_S\neq 1$.
That is why CTF matching fails for rippled graphene~\cite{Yu2017} and
why stage~1 is biased.

If the object is known to be a single sinusoid, the unknown is
$(\varphi,x_0)$ (period $\Lambda$ is read from the fundamental of $I$).
The aperture keeps
\begin{equation}
M=2n_{\mathrm{ap}}+1,
\qquad
n_{\mathrm{ap}}=\bigl\lfloor \qap\Lambda\bigr\rfloor
\end{equation}
modes (Lean: \texttt{nAperture}, \texttt{card\_modeSet}).
The FO image is a double sum of $M^2$ terms
(\texttt{sinusoidalFOImage}, \texttt{card\_modePairs}).
Least squares on $\varphi$ with analytic Jacobian
$J_n'=\tfrac12(J_{n-1}-J_{n+1})$ is $O(k M^2)$, $k\sim 5$--$20$.
Dropped modes obey the $O(1/n^2)$ bound already proved
(\texttt{besselJ\_bound}, \texttt{besselJ\_outside\_modeSet}).

This reduction does \emph{not} apply to a general 2D micrograph:
there is no Bessel ladder.
Then stage~2 is Gauss--Newton in the aperture basis, initialized by
stage~1.

\paragraph{Skip test.}
Compute the FO residual of the stage~1 reconstruction.
If $r_1^2/\sigma^2$ is consistent with noise, quadratic products are
undetectable and stage~2 cannot help.
Do not skip merely because $\Gamma_C\Gamma_S\approx 1$ on the axes:
the linear inverse never saw off-axis envelopes.

\section{What is proved, and what is not}
\label{sec:lean}

\begin{center}
\small
\begin{tabular}{@{}ll@{}}
\toprule
Claim & Lean \\
\midrule
Tikhonov unique min, closed form & \texttt{LeemFO/Tikhonov.lean} \\
Bias--noise identity and sharp triangle bound & same \\
$O(KN\log N)$ cost model vs dense $KN^2$ & \texttt{reconstructCost\_*} \\
Aperture support of $\RFO(\cdot,0)$ & \texttt{LeemFO/LinearInverse.lean} \\
Gauge, vacuum Jacobian, quadratic remainder & same \\
$K=1$ not injective; NAC CTF zeros & same \\
Mode count $M=2\lfloor\qap\Lambda\rfloor+1$ & \texttt{LeemFO/Inverse.lean} \\
Jacobi--Anger and Bessel decay & \texttt{LeemFO/PhaseObject.lean} \\
\midrule
FFT exists / equals the continuum FT & informal \\
Gaussian/Poisson noise, Cram\'er--Rao & informal \\
GN local quadratic convergence & informal \\
Image alignment, current-to-defocus map & informal \\
``Fastest among all conceivable algorithms'' & only DFT vs dense bilinear \\
\bottomrule
\end{tabular}
\end{center}

\section{Conclusion}
\label{sec:end}

For a 2D experimental LEEM focal series the fastest method that still
reconstructs a scientifically usable object --- filling CTF zeros,
using the FO/CTF slice already formalized in this repository, with a
proved unique regularized inverse and a proved error identity --- is
\textbf{multi-defocus Fourier-diagonal Tikhonov} (Schiske/Wiener,
$O(KN\log N)$).
That is stage~1.
Stage~2, Gauss--Newton on bilinear FO, is the FO-faithful refinement
for strong-phase topography and is unnecessary when the stage~1 FO
residual is at the noise floor.
For the 1D sinusoid that Yu~\emph{et al.}\ use to exhibit CTF failure,
stage~2 collapses to Jacobi--Anger fitting and is then the fastest
still-correct inverse of all.

This is the inverse of the forward model in~\cite{Yu2019}, not a
replacement for it: Yu forward-matches; Duden inverts linearly;
the two-stage pipeline does both, with Lean certificates for the
linear stage.

\begin{thebibliography}{9}
\bibitem{Yu2019}
K.~M. Yu, K.~L.~W. Lau, M.~S. Altman,
Ultramicroscopy \textbf{200} (2019) 160--168.

\bibitem{Yu2017}
K.~M. Yu, A. Locatelli, M.~S. Altman,
Ultramicroscopy \textbf{183} (2017) 109--116.

\bibitem{Schramm2012}
S.~M. Schramm, A.~B. Pang, M.~S. Altman, R.~M. Tromp,
Ultramicroscopy \textbf{115} (2012) 88--108.

\bibitem{Duden2014}
T.~Duden, A.~Thust, C.~Kumpf, F.~S. Tautz,
Microsc.\ Microanal.\ \textbf{20} (2014) 968--973.

\bibitem{Coene1996}
W.~M.~J. Coene, A.~Thust, M.~Op de Beeck, D.~Van Dyck,
Ultramicroscopy \textbf{64} (1996) 109--135.

\bibitem{Teague1983}
M.~R. Teague,
J.\ Opt.\ Soc.\ Am.\ \textbf{73} (1983) 1434--1441.

\bibitem{Paganin}
D.~Paganin and K.~A. Nugent,
Phys.\ Rev.\ Lett.\ \textbf{80} (1998) 2586--2589.

\bibitem{Schiske1968}
P.~Schiske,
in \emph{Image Processing and Computer-Aided Design in Electron Optics},
ed.\ P.~W. Hawkes (Academic Press, 1973); originally 1968.
\end{thebibliography}

\end{document}
